\section{Minimal combined skeleton}

We can now write the partial combination without claiming a complete HLS theory.

Let \(F_t\) denote the distribution of incoming problems and \(S_t\) the current
knowledge/competence state. The operational organization is
\begin{equation}
a_t=\mathcal R(F_t,S_t,\kappa_t).
\end{equation}
Executing this organization generates experience
\begin{equation}
x_t=\mathcal X(a_t,F_t,\omega_t),
\end{equation}
where \(\omega_t\) denotes execution outcomes or other factors determining realized
experience. Competence/knowledge then evolves according to
\begin{equation}
S_{t+1}=\mathcal L(S_t,x_t,\theta_t).
\end{equation}
The next period therefore uses
\begin{equation}
a_{t+1}=\mathcal R(F_{t+1},S_{t+1},\kappa_{t+1}).
\end{equation}

Combining the equations gives the feedback system
\begin{equation}
S_{t+1}
=
\mathcal L\!\left(
S_t,
\mathcal X\!\left(
\mathcal R(F_t,S_t,\kappa_t),
F_t,\omega_t
\right),
\theta_t
\right).
\label{eq:combined}
\end{equation}

Equation~\eqref{eq:combined} is the central mathematical statement of this note. It should be
read as a structural composition, not as a solved model. Each operator still requires a
specific instantiation before optimization or theorem development is possible.

\subsection{What the combination adds}

Neither beam alone contains the full feedback.

Beam~1 can organize current knowledge efficiently but, in its base form, does not let the
organization's own operation change future competence.

Beam~2 allows current allocation to change future competence but does not supply Garicano's
problem-matching and escalation structure.

Their combination introduces the possibility that an operationally sensible allocation at
time \(t\) changes the competence state used to organize work at \(t+1\). This is the precise
sense in which present organization and future capability become linked.

No stronger claim is needed at this stage.
